\mysection{8086メモリモデル}
\myindex{Intel!8086!Memory model}
\myindex{MS-DOS}
\label{8086_memory_model}

MS-DOS用の16ビットプログラムやWin16
（\myref{dongle_16bit_dos}または\myref{win16_near_far_pointers}）を扱う場合、
ポインターが2つの16ビット値で構成されていることがわかります。
これらは何を意味するのでしょうか？そう、これはMS-DOSと8086のもう一つの奇妙な遺物です。

8086/8088は16ビットCPUでしたが、RAMの20ビットアドレスをアドレス指定できました
（したがって1MBの外部メモリにアクセスできました）。

外部メモリアドレス空間は\ac{RAM}（最大640KB）、
\ac{ROM}、ビデオメモリ用ウィンドウ、EMSカード、\etc{}に分割されていました。

8086/8088が実際には8ビット8080 CPUの継承者であったことも思い出しましょう。

8080は16ビットメモリ空間を持ち、つまり64KBしかアドレス指定できませんでした。

そして、おそらく古いソフトウェア移植の理由で\footnote{著者はここで100％確信していません}、
8086は1MBアドレス空間内に配置された多くの64KBウィンドウを同時にサポートできます。

これは玩具レベルの仮想化の一種です。
\myindex{x86!\Registers!CS}
\myindex{x86!\Registers!DS}
\myindex{x86!\Registers!ES}
\myindex{x86!\Registers!SS}

すべての8086レジスタは16ビットなので、より多くをアドレス指定するために、特別なセグメントレジスタ（CS、DS、ES、SS）が導入されました。

各20ビットポインターは、セグメントレジスタとアドレスレジスタペア（例：DS:BX）の値を使用して次のように計算されます：

\begin{center}
$real\_address = (segment\_register \ll 4) + address\_register$
\end{center}

例えば、古いIBM PC互換機のグラフィックス（\ac{EGA}、\ac{VGA}）ビデオ\ac{RAM}ウィンドウのサイズは64KBです。

これにアクセスするには、0xA000の値をセグメントレジスタの一つ（例：DS）に格納する必要があります。

そうするとDS:0はビデオ\ac{RAM}の最初のバイトをアドレス指定し、DS:0xFFFF〜はRAMの最後のバイトをアドレス指定します。

ただし、20ビットアドレスバス上の実際のアドレスは0xA0000から0xAFFFFの範囲になります。

プログラムには0x1234のようなハードコードされたアドレスが含まれている可能性がありますが、\ac{OS}は任意のアドレスでプログラムをロードする必要がある場合があるため、プログラムがRAM内のどこに配置されているかを気にする必要がないように、セグメントレジスタ値を再計算します。

したがって、古いMS-DOS環境のポインターは実際にはセグメントアドレスとセグメント内のアドレス、つまり2つの16ビット値で構成されていました。20ビットで十分でしたが、アドレスを非常に頻繁に再計算する必要がありました：スタックにより多くの情報を渡すことは、より良い空間/利便性のバランスのように見えました。

ちなみに、このすべてのために、64KBより大きなメモリブロックを割り当てることはできませんでした。

\myindex{Intel!80286}
\myindex{Intel!80386}

セグメントレジスタは80286でセレクタとして再利用され、異なる機能を果たしました。

\myindex{MS-DOS!DOS extenders}

80386 CPUとより大きな\ac{RAM}を持つコンピューターが導入されたとき、
MS-DOSはまだ人気があったため、DOSエクステンダーが登場しました：これらは実際には
「真面目な」\ac{OS}への一歩であり、
CPUをプロテクトモードに切り替え、MS-DOS下で実行する必要があるプログラムに対してはるかに優れたメモリ\ac{API}を提供しました。

広く人気のある例には、DOS/4GW（DOOMビデオゲームがこれ用にコンパイルされました）、Phar Lap、PMODEがあります。
\par
\myindex{Windows!Windows 3.x}
\myindex{Windows!Win32}

ちなみに、同じメモリアドレス指定方法がWin32以前のWindows 3.xの16ビットラインでも使用されていました。